% Final layout: two selected boundary rays, matching the edited reference.
% Direct TikZ -> XDV -> SVG. Canvas: exactly 8.8 x 5.4 cm.
\def\pgfsysdriver{pgfsys-dvisvgm.def}
\documentclass[tikz,border=0pt]{standalone}
\usepackage{fontspec,amsmath,bm}
\setsansfont{Arial}
\usetikzlibrary{arrows.meta,calc}
\definecolor{ink}{HTML}{243746}
\definecolor{beam}{HTML}{FF7F0E} % Matplotlib C1
\definecolor{optic}{HTML}{858585}
\definecolor{signal}{HTML}{2CA02C} % Matplotlib C2
% Physical inputs: E in keV, Ge lattice parameter in Angstrom, R in cm.
\def\Energy{25.514}
\def\Lattice{5.6579}
\def\Radius{20.6}
% O-2094: aluminum substrate 52 x 27 mm, center thickness 10 mm.
% 52 mm curved section shown symbolically in this 2D optical layout;
% the 27 mm cylinder-axis dimension is out of the icon plane.
\def\SupportWidth{5.2}
\def\SupportDepth{2.7}
\def\SupportThickness{1.0}
% Visual exaggeration only: leave Radius and the nominal arms unchanged.
\def\CurvatureExaggeration{3.0}
\pgfmathsetmacro\Spacing{\Lattice/sqrt(128)}
\pgfmathsetmacro\BraggAngle{asin(12.39841984332/(2*\Spacing*\Energy))}
\pgfmathsetmacro\Arm{\Radius/sin(\BraggAngle)}
% Drawing scale: paper cm / physical cm. Same scale on x and y.
\pgfmathsetmacro\Scale{3.45/\Arm}
\newcommand{\num}[1]{\pgfmathprintnumber[fixed,precision=2]{#1}}
\begin{document}\boldmath
\begin{tikzpicture}[x=1cm,y=1cm,line width=.3pt,line cap=round,line join=round,
 every node/.style={font=\sffamily\fontsize{6.5}{7.5}\selectfont,text=ink,inner sep=1pt},
 ray/.style={draw=beam,line width=.40pt,-{Stealth[length=.85mm,width=.55mm]}},
 guide/.style={draw=ink!45,dash pattern=on 1.5pt off 1.5pt,line width=.25pt}]
\path[use as bounding box] (0,0) rectangle (8.8,5.4);
\begin{pgfinterruptboundingbox}
\fill[white] (0,0) rectangle (8.8,5.4);
% Isotropic layout enlargement; no anisotropic stretching of rays or text.
% Center the optical construction in the wider canvas.
\begin{scope}[shift={(.62,0)},scale=1.08]



\coordinate (S) at (.55,3.65);
\coordinate (C) at (4,3.65);
\coordinate (F) at ($(C)+({3.45*cos(2*\BraggAngle)},{-3.45*sin(2*\BraggAngle)})$);
% Two illustrative boundary rays, selected by both slit apertures.
% Their input directions intersect at a virtual source V between the slits.
% Envelope in this view is conceptual: sagittal focusing is out of plane.
\def\VirtualSourceX{1.80}
\def\CrystalIntercept{1.90} % physical cm, within the 2.6 cm substrate half-width
\pgfmathsetmacro\EdgeSag{\CurvatureExaggeration*(sqrt(\Radius*\Radius-\CrystalIntercept*\CrystalIntercept)-\Radius)}
\pgfmathsetmacro\UpperCX{4+\Scale*(-\CrystalIntercept*cos(\BraggAngle)+\EdgeSag*sin(\BraggAngle))}
\pgfmathsetmacro\UpperCY{3.65+\Scale*(\CrystalIntercept*sin(\BraggAngle)+\EdgeSag*cos(\BraggAngle))}
\pgfmathsetmacro\LowerCX{4+\Scale*(\CrystalIntercept*cos(\BraggAngle)+\EdgeSag*sin(\BraggAngle))}
\pgfmathsetmacro\LowerCY{3.65+\Scale*(-\CrystalIntercept*sin(\BraggAngle)+\EdgeSag*cos(\BraggAngle))}
\pgfmathsetmacro\SlopeUpper{(\UpperCY-3.65)/(\UpperCX-\VirtualSourceX)}
\pgfmathsetmacro\SlopeLower{(\LowerCY-3.65)/(\LowerCX-\VirtualSourceX)}
\coordinate (Cupper) at (\UpperCX,\UpperCY);
\coordinate (Clower) at (\LowerCX,\LowerCY);
\coordinate (Supper) at (.55,{3.65+\SlopeUpper*(.55-\VirtualSourceX)});
\coordinate (Slower) at (.55,{3.65+\SlopeLower*(.55-\VirtualSourceX)});
\draw[ray] (Supper)--(Cupper);
\draw[ray] (Slower)--(Clower);
\draw[ray] (Cupper)--(F);
\draw[ray] (Clower)--(F);
% p and q label the nominal central-arm lengths, not the envelope ray lengths.
\node[above] at (2.15,3.76) {${p}$};
\node[rotate={-2*\BraggAngle},above] at ($(C)!.57!(F)+(-0.098808,0.522709)$) {${q}$};
% Tube source; emission point S fixes the origin of the optical path.
\draw[fill=ink!7,draw=ink] (.22,3.38) rectangle (.49,3.92);
\draw[fill=ink!16,draw=ink] (.49,3.53) rectangle (.55,3.77);
\fill[beam] (Supper) circle (.017);
\fill[beam] (Slower) circle (.017);
\node[align=center] at (0.294635,4.281802) {X-Ray Tube\\Source};
% Apertures enclose both line intersections, with a stroke-width clearance.
% Their values are illustrative and must not be read as design slit settings.
\foreach \x/\lab in {1.15/1,3.12/2}{
 \pgfmathsetmacro\RayUpper{3.65+max(\SlopeUpper*(\x-\VirtualSourceX),\SlopeLower*(\x-\VirtualSourceX))}
 \pgfmathsetmacro\RayLower{3.65+min(\SlopeUpper*(\x-\VirtualSourceX),\SlopeLower*(\x-\VirtualSourceX))}
 \draw[line width=.7pt,ink]
  (\x,3.32)--(\x,{\RayLower-.025})
  (\x,{\RayUpper+.025})--(\x,3.98);

}
\node[anchor=base] at (1.197533,3.02) {Entrance Slit};
\node[anchor=base] at (3.158051,3.02) {Collimating Slit};
% O-2094 chord length and flat back; front-face curvature visually exaggerated.
% Curvature direction remains symbolic in this dispersive-plane scheme.
\begin{scope}[shift={(C)},rotate={-\BraggAngle}]
 \begin{scope}[x={\Scale cm},y={\Scale cm}]
  \draw[fill=optic!35,draw=optic!80!black,line width=.25pt]
   plot[domain={-\SupportWidth/2}:{\SupportWidth/2},samples=41,variable=\u]
    (\u,{\CurvatureExaggeration*(sqrt(\Radius*\Radius-\u*\u)-\Radius)})
   -- ({\SupportWidth/2},\SupportThickness)
   -- ({-\SupportWidth/2},\SupportThickness) -- cycle;
 \end{scope}
 \draw[guide] (-1,0)--({-\SupportWidth*\Scale/2},0);
\end{scope}
\node at (4.507800,3.912956) {Ge(880)};



% Bragg angle measured against the nominal incident direction.
\draw[guide] (3.30,3.65)--(C);
\draw[ink] ($(C)+(180:.68)$) arc[start angle=180,end angle={180-\BraggAngle},radius=.68];
\node at (3.001412,4.354912) {${\theta_B=\num{\BraggAngle}^{\circ}}$};
% Focus at the sample, not at the fluorescence detector.
% Sample surface is exactly 45 degrees to the incident propagation direction.
\pgfmathsetmacro\SampleAngle{-2*\BraggAngle+45}
\begin{scope}[shift={(F)},rotate=\SampleAngle]
 \draw[draw=optic!80!black,fill=optic!35] (-.30,-.035) rectangle (.30,.035);
 \draw[guide] (-.62,0)--(-.30,0);
\end{scope}
\draw[guide] (F)--++({180-2*\BraggAngle}:.48);
\draw[ink!70] ($(F)+({180-2*\BraggAngle}:.44)$)
 arc[start angle={180-2*\BraggAngle},end angle={180+\SampleAngle},radius=.44];
\node at ($(F)+(-.67,.39)$) {$45^{\circ}$};
\fill[beam] (F) circle (.025);
\node at (5.336047,0.489682) {Ag Sample};
% Illustrative fluorescence collection direction; not a constrained distance.
\coordinate (D) at ($(F)+(.62,1.00)$);
% Three fluorescence arrows fan out toward the CZT entrance face.
\foreach \offset in {-.16,0,.16}{
 \draw[draw=signal,line width=.40pt,shorten <=.6mm,shorten >=.15mm,
  -{Stealth[length=.8mm,width=.5mm]}]
  (F)--($(D)+({\offset*cos(32)},{-\offset*sin(32)})$);
}
\begin{scope}[shift={(D)},rotate=-32]
 \draw[fill=optic!35,draw=optic!80!black] (-.22,0) rectangle (.22,.22);
\end{scope}
\node at (6.564949,2.303608) {Multichannel CZT};
% Cartesian directions bottom left; physical scale bottom right.
\draw[-{Stealth[length=1mm]},ink] (-.30,.42)--(.30,.42) node[right] {${x}$};
\draw[-{Stealth[length=1mm]},ink] (-.30,.42)--(-.30,1.02) node[above] {${y}$};
\coordinate (ScaleStart) at (6.38,.20);
\draw[ink,line width=.45pt] (ScaleStart)--++({10*\Scale},0);
\draw[ink] ($(ScaleStart)+(0,-.06)$)--++(0,.12)
 ($(ScaleStart)+({10*\Scale},-.06)$)--++(0,.12);
\node[anchor=south] at ($(ScaleStart)+({5*\Scale},.08)$) {10 cm};
\node at (2.212039,1.775024) {${p=q=\dfrac{R}{\sin\theta_B}=\num{\Arm}\,\mathrm{cm}}$};
\end{scope}
\end{pgfinterruptboundingbox}
\end{tikzpicture}
\end{document}
